%------------------------------------------------------------------------------
% CHAPTER 3: SPECIFIC REQUIREMENTS
%------------------------------------------------------------------------------
\chapter{Specific Requirements}
\label{ch:requirements}

\section{External Interface Requirements}
\label{sec:interfaces}

\subsection{User Interfaces}

\begin{itemize}
    \item \textbf{Sign Up/Login Page}: The system must include a Sign Up/Login page that allows Users to register and authenticate their credentials for accessing the platform. This functionality ensures secure access for registered users.

    \item \textbf{Home Page and Search Bar}: The system must include a home page which has in the bottom buttons to access personal dashboard and profile management and in the top should have a search bar to insert origin and destination points of a trip.

    \item \textbf{Map interface}: When Users search a path, a map interface representing the possible paths, ordered by score, should be presented to the user. Users can tap on various paths and look at some details.

    \item \textbf{Dashboard}: The system must include a dashboard section where registered users can look at all their past registered trips and related statistics (e.g., total distance, trip time, average speed).
\end{itemize}

\subsection{Hardware Interfaces}

The Best Bike Paths (BBP) application interacts directly with the following hardware components hosted on the user's device:

\begin{itemize}
    \item \textbf{GPS Receiver}: The system shall interface with the GPS receiver to obtain real-time geolocation coordinates for trip tracking and navigation.

    \item \textbf{Inertial Sensors}:
    \begin{itemize}
        \item \textbf{Accelerometer}: Used to measure acceleration to detect sudden jolts indicative of potholes.
        \item \textbf{Gyroscope}: Used to measure orientation and angular velocity to distinguish between road anomalies and phone handling movements.
    \end{itemize}

    \item \textbf{Touchscreen Display}: Used for all user inputs and visual output.
\end{itemize}

\subsection{Software Interfaces}

The system interacts with the following external software services and APIs:

\begin{itemize}
    \item \textbf{Map Rendering Service}: The system shall interface with an external mapping provider (e.g., Google Maps Platform, OpenStreetMaps) to retrieve map tiles and perform geocoding services.

    \item \textbf{Meteorological Data Service}: The system shall communicate with a third-party weather API to fetch real-time environmental data (or past for registered trips) based on the user's location.
\end{itemize}

\subsection{Communication Interfaces}

The application can be accessed through a stable internet connection to perform all actions and access all functionalities. Connection is not required while biking since local buffering for trip recording is present.

\section{Functional Requirements}
\label{sec:functional}

\subsection{Sign up and log in}

\begin{itemize}
    \item \textbf{[R1]}: The system must allow an unregistered User to sign up
    \item \textbf{[R2]}: The system must allow registered Users to login
\end{itemize}

\subsection{Manual Trip Tracking}

\begin{itemize}
    \item \textbf{[R3]}: The system must allow a Registered User to manually insert a bike path by listing the name and the status of each street comprising the path.
    \item \textbf{[R4]}: The system must store manually inserted paths in the User's private trip history.
\end{itemize}

\subsection{Automatic Trip Tracking}

\begin{itemize}
    \item \textbf{[R5]}: The system must allow a Registered User to automatically record a trip, using device sensors to detect when he starts and stops biking.
    \item \textbf{[R6]}: The system must automatically identify the start of a Registered User's biking activity and initiate the recording process.
    \item \textbf{[R7]}: The system must automatically start GPS-based tracking when biking is detected.
    \item \textbf{[R8]}: The system must collect GPS coordinates to reconstruct the path followed by the User.
    \item \textbf{[R9]}: The system must collect accelerometer and gyroscope data to detect potential anomalies (e.g., potholes).
    \item \textbf{[R10]}: The system must store complete trip statistics (distance, average speed, duration and, when available, weather conditions at the time of recording).
    \item \textbf{[R11]}: The system must allow a registered user to view the complete statistics (distance, average speed, duration and, when available, weather conditions) related to his automated-recorded trip.
    \item \textbf{[R12]}: The system must buffer trip data locally when network connectivity is unavailable and must synchronize buffered trip data once connectivity is restored.
    \item \textbf{[R13]}: The system must detect when biking ends and automatically stop tracking.
    \item \textbf{[R14]}: In case of anomalies detection, the system must allow a Registered User to confirm or reject each detected anomaly within 72 hours from the trip recording.
    \item \textbf{[R15]}: The system must allow a Registered User to publish a validated trip.
\end{itemize}

\subsection{Trip Storage and Personal Analytics}

\begin{itemize}
    \item \textbf{[R16]}: The system must allow Registered Users to store all their recorded trips (manually or automatically).
    \item \textbf{[R17]}: The system must allow a Registered User to view his private trip history, including relatively collected information.
    \item \textbf{[R18]}: The system must allow a Registered User to review the statistics related to a bike path, such as average speed, and other performance metrics.
\end{itemize}

\subsection{Path Search, Routing and Visualization}

\begin{itemize}
    \item \textbf{[R19]}: The system must allow any user to access published paths
    \item \textbf{[R20]}: The system must allow any User to search the bike path(s) between a given origin and a given destination.
    \item \textbf{[R21]}: The system must visualize the paths on a map interface.
    \item \textbf{[R22]}: The system must compute a Path Score based on path status and estimated trip time.
    \item \textbf{[R23]}: The system must rank available paths based on their Path Score.
    \item \textbf{[R24]}: The system must allow a User to inspect public details of a selected path, such as status, total distance, trip time, and validated anomalies.
    \item \textbf{[R25]}: The system shall query an external weather service to retrieve current weather conditions during the path search and display relevant warnings to the user.
\end{itemize}

\subsection{Merging \& Freshness-Based Status Aggregation}

\begin{itemize}
    \item \textbf{[R26]}: The system shall append historical weather data (temperature, wind speed and direction) to the trip logs of registered users for performance analysis.
    \item \textbf{[R27]}: The system must merge publishable information from multiple Users referring to the same path, considering their freshness and the number of confirming contributions.
\end{itemize}

\section{Use Cases}
\label{sec:use_cases}

\subsection{Use Case Diagram}

\begin{figure}[H]
    \centering
    % \includegraphics[width=0.9\textwidth]{use_case_diagram.png}
    \caption{Use Case Diagram}
    \label{fig:use_case_diagram}
\end{figure}

\subsection{Use Case Descriptions}

\subsubsection{[UC1] - Sign Up}

\begin{table}[H]
\centering
\begin{tabularx}{\textwidth}{|l|X|}
\hline
\textbf{Participating Actors} & Guest User \\ \hline
\textbf{Entry Conditions} & The Guest User is not authenticated and wants to become a Registered User \\ \hline
\textbf{Flow of events} &
\begin{enumerate}
    \item The Guest User clicks on "sign up"
    \item BBP displays the registration form
    \item The User inputs email, username, password
    \item The User clicks "create account"
    \item BBP validates the input syntax
    \item BBP creates a new profile in the database
    \item BBP confirms creation and redirects to the Login Page
\end{enumerate} \\ \hline
\textbf{Exit Conditions} & The User has successfully created a new account. \\ \hline
\textbf{Exceptions} &
E-mail already in use: BBP notifies the User and suggests logging in

E-mail not valid: BBP notifies the User and suggest correcting it \\ \hline
\end{tabularx}
\caption{Use Case 1: Sign Up}
\label{tab:uc1}
\end{table}

\begin{figure}[H]
    \centering
    % \includegraphics[width=0.7\textwidth]{uc1_sequence.png}
    \caption{Sequence Diagram for UC1: Sign Up}
    \label{fig:uc1_sequence}
\end{figure}

\subsubsection{[UC2] - Log In}

\begin{table}[H]
\centering
\begin{tabularx}{\textwidth}{|l|X|}
\hline
\textbf{Participating Actors} & Registered User \\ \hline
\textbf{Entry Conditions} & The Registered User is not currently authenticated \\ \hline
\textbf{Flow of events} &
\begin{enumerate}
    \item The User clicks on "login"
    \item BBP displays the login form
    \item The User inputs their credentials (email/username, password)
    \item The User clicks the "login" button
    \item BBP hashes the password and validates it against the database
    \item BBP retrieves the user's profile and starts a session
    \item BBP redirects the user to the Home Page
\end{enumerate} \\ \hline
\textbf{Exit Conditions} & The User is authenticated and can access restricted features. \\ \hline
\textbf{Exceptions} & Invalid Credentials: BBP displays an error message and invites the user to retry. \\ \hline
\end{tabularx}
\caption{Use Case 2: Log In}
\label{tab:uc2}
\end{table}

\begin{figure}[H]
    \centering
    % \includegraphics[width=0.7\textwidth]{uc2_sequence.png}
    \caption{Sequence Diagram for UC2: Log In}
    \label{fig:uc2_sequence}
\end{figure}

\subsubsection{[UC3] - Record Trip Automatically}

\begin{table}[H]
\centering
\begin{tabularx}{\textwidth}{|l|X|}
\hline
\textbf{Participating Actors} & Registered User, Mobile Sensors, Weather Service \\ \hline
\textbf{Entry Conditions} & The Registered User is ready to automatically record his biking activity \\ \hline
\textbf{Flow of events} &
\begin{enumerate}
    \item User clicks on "Start automatic recording"
    \item BBP enters detection mode and starts monitoring GPS speed
    \item BBP detects biking activity (speed > threshold)
    \item BBP starts the recording session
    \item BBP continuously logs GPS coordinates and timestamps (trajectory)
    \item BBP monitors accelerometer/gyroscope for anomalies (e.g., potholes)
    \item [Alt: speed < threshold] BBP pauses data acquisition and returns to step 2
    \item User presses "End Automatic Record" button
    \item BBP stores the detailed raw data (GPS, trajectory, timestamps)
    \item BBP retrieves the historical meteorological data from the Weather Service and link it with the trip
    \item BBP analyzes sensors data to identify potholes and stores them as PENDING
    \item BBP calculates the specific statistics (total distance, average speed, and duration) and stores them
    \item BBP saves the trip locally
    \item BBP synchronizes the trip with the server
\end{enumerate} \\ \hline
\textbf{Exit Conditions} & A new trip is recorded with its potholes, statistics and weather conditions \\ \hline
\textbf{Exceptions} & Connectivity Loss: BBP buffers data locally until connection is restored. \\ \hline
\end{tabularx}
\caption{Use Case 3: Record Trip Automatically}
\label{tab:uc3}
\end{table}

\begin{figure}[H]
    \centering
    % \includegraphics[width=\textwidth]{uc3_sequence.png}
    \caption{Sequence Diagram for UC3: Record Trip Automatically}
    \label{fig:uc3_sequence}
\end{figure}

\subsubsection{[UC4] - Record Trip Manually}

\begin{table}[H]
\centering
\begin{tabularx}{\textwidth}{|l|X|}
\hline
\textbf{Participating Actors} & Registered User, Map Service \\ \hline
\textbf{Entry Conditions} & The Registered User is ready to manually insert trip information \\ \hline
\textbf{Flow of events} &
\begin{enumerate}
    \item The user selects "Manual Record"
    \item BBP presents a form to insert a street name
    \item [Loop] The user specifies a street name
    \begin{enumerate}
        \item The BBP queries the Map Service with provided street name
        \item The BBP displays found streets to the user
        \item Registered User selects one of the found streets
        \item [End Loop] BBP add the chosen street to the trip
    \end{enumerate}
    \item The user submits the form
    \item BBP confirms the submission
    \item BBP saves the data
\end{enumerate} \\ \hline
\textbf{Exit Conditions} & A new manual trip is recorded \\ \hline
\textbf{Exceptions} &
Invalid Location: The system cannot identify the street provided. abort

Map Service is unavailable: BBP shows an error regarding the map service and cancel the operation suggesting trying again later \\ \hline
\end{tabularx}
\caption{Use Case 4: Record Trip Manually}
\label{tab:uc4}
\end{table}

\begin{figure}[H]
    \centering
    % \includegraphics[width=\textwidth]{uc4_sequence.png}
    \caption{Sequence Diagram for UC4: Record Trip Manually}
    \label{fig:uc4_sequence}
\end{figure}

\subsubsection{[UC5] - Validate Obstacle}

\begin{table}[H]
\centering
\begin{tabularx}{\textwidth}{|l|X|}
\hline
\textbf{Participating Actors} & Registered User, Map Service \\ \hline
\textbf{Entry Conditions} & The registered user is on its personal history page and is ready to validate an automatically recorded Trip \\ \hline
\textbf{Flow of events} &
\begin{enumerate}
    \item The Registered User chooses one of its automatically recorded private trips
    \item BBP highlights potential selected trip's anomalies on the map
    \item The user reviews a specific anomaly
    \item [Alt] The user confirms the detection
    \begin{enumerate}
        \item [Alt] The BBP updates pothole status to CONFIRMED
    \end{enumerate}
    \item [Alt2] The user rejects the detection
    \begin{enumerate}
        \item [Alt2] The BBP updates pothole status to REJECTED
    \end{enumerate}
    \item BBP updates the trip accordingly
    \item The user saves the validated trip
\end{enumerate} \\ \hline
\textbf{Exit Conditions} & Trip obstacles are validated and trip is ready to for optional publication \\ \hline
\textbf{Exceptions} & \\ \hline
\end{tabularx}
\caption{Use Case 5: Validate Obstacle}
\label{tab:uc5}
\end{table}

\begin{figure}[H]
    \centering
    % \includegraphics[width=\textwidth]{uc5_sequence.png}
    \caption{Sequence Diagram for UC5: Validate Obstacle}
    \label{fig:uc5_sequence}
\end{figure}

\subsubsection{[UC6] - Publish Trip}

\begin{table}[H]
\centering
\begin{tabularx}{\textwidth}{|l|X|}
\hline
\textbf{Participating Actors} & Registered User \\ \hline
\textbf{Entry Conditions} & The registered user has at least one private trip ready to publish \\ \hline
\textbf{Flow of events} &
\begin{enumerate}
    \item The user selects a private trip from his history
    \item The user clicks "publish"
    \item BBP changes the trip status to "public"
    \item BBP anonymizes trip data
    \item BBP triggers the Data Merging algorithm to integrate anonymized data with already existing data (considering freshness/confirmations)
    \item BBP updates the global path score
    \item BBP confirms publication
\end{enumerate} \\ \hline
\textbf{Exit Conditions} & Trip data has been integrated in public path and it's available to all users \\ \hline
\textbf{Exceptions} & \\ \hline
\end{tabularx}
\caption{Use Case 6: Publish Trip}
\label{tab:uc6}
\end{table}

\begin{figure}[H]
    \centering
    % \includegraphics[width=\textwidth]{uc6_sequence.png}
    \caption{Sequence Diagram for UC6: Publish Trip}
    \label{fig:uc6_sequence}
\end{figure}

\subsubsection{[UC7] - View Private Trip}

\begin{table}[H]
\centering
\begin{tabularx}{\textwidth}{|l|X|}
\hline
\textbf{Participating Actors} & Registered User, Map Service \\ \hline
\textbf{Entry Conditions} & The registered user has recorded at least one trip and wants to view it in details \\ \hline
\textbf{Flow of events} &
\begin{enumerate}
    \item The user clicks on "My Trips"
    \item BBP displays the list of the Registered User trips over time
    \item The user selects a specific trip from his history list
    \item BBP fetches map data about chosen trip's coordinates
    \item BBP displays a detailed view showing trip trajectory on the map, computed statistics, found/confirmed potholes and weather conditions faced during the ride
\end{enumerate} \\ \hline
\textbf{Exit Conditions} & The user successfully visualizes the detailed trip performance and context of a single specific trip \\ \hline
\textbf{Exceptions} & Map Service is unavailable: BBP shows an error regarding the map but still shows all other things \\ \hline
\end{tabularx}
\caption{Use Case 7: View Private Trip}
\label{tab:uc7}
\end{table}

\begin{figure}[H]
    \centering
    % \includegraphics[width=\textwidth]{uc7_sequence.png}
    \caption{Sequence Diagram for UC7: View Private Trip}
    \label{fig:uc7_sequence}
\end{figure}

\subsubsection{[UC8] - Delete Private Trip}

\begin{table}[H]
\centering
\begin{tabularx}{\textwidth}{|l|X|}
\hline
\textbf{Participating Actors} & Registered User \\ \hline
\textbf{Entry Conditions} & The registered user has recorded at least one trip and wants to delete it \\ \hline
\textbf{Flow of events} &
\begin{enumerate}
    \item The registered user clicks on "My Trips"
    \item BBP displays the list of the user trips over time
    \item The user keeps pressed on a specific trip from his history list
    \item BBP shows an additional in-context menu
    \item The user selects the "Delete" option
    \item BBP ask for confirmation
    \item [Alt] The user confirm the operation
    \begin{enumerate}
        \item [Alt] BBP deletes the selected Trip from User history
    \end{enumerate}
    \item [Alt] The user cancels the operation
    \begin{enumerate}
        \item [Alt] BBP aborts the operation and displays the user trips history
    \end{enumerate}
\end{enumerate} \\ \hline
\textbf{Exit Conditions} & Trip is removed from the registered user's private history. If previously published, anonymized data remains in aggregated PATH information. \\ \hline
\textbf{Exceptions} &
Invalid Location: The system cannot identify the street provided. abort

Map Service is unavailable: BBP shows an error regarding the map service and cancel the operation suggesting trying again later \\ \hline
\end{tabularx}
\caption{Use Case 8: Delete Private Trip}
\label{tab:uc8}
\end{table}

\begin{figure}[H]
    \centering
    % \includegraphics[width=\textwidth]{uc8_sequence.png}
    \caption{Sequence Diagram for UC8: Delete Private Trip}
    \label{fig:uc8_sequence}
\end{figure}

\subsubsection{[UC9] - Search Path}

\begin{table}[H]
\centering
\begin{tabularx}{\textwidth}{|l|X|}
\hline
\textbf{Participating Actors} & Any User (Guest/Registered), Map Service \\ \hline
\textbf{Entry Conditions} & The User is ready to search paths between two points \\ \hline
\textbf{Flow of events} &
\begin{enumerate}
    \item The user clicks on the search bar
    \item BBP displays search form
    \item [Loop for Origin and Destination endpoints]
    \begin{enumerate}
        \item The user enters endpoint street name
        \item BBP queries the Map Service with provided street name
        \item BBP displays found streets to the user
        \item The User selects one of the found streets
        \item BBP add selected street as endpoint value
    \end{enumerate}
    \item [End Loop]
    \item BBP searches paths with provided endpoints
    \item BBP ranks the found paths by previously computed/updated score
    \item BBP shows the ranked list of paths to the user
\end{enumerate} \\ \hline
\textbf{Exit Conditions} & User is shown of available paths ready to be viewed in details \\ \hline
\textbf{Exceptions} &
Invalid Location: The system cannot identify the street provided. abort

Map Service is unavailable: BBP shows an error regarding the map service and cancel the operation suggesting trying again later

No Paths Found: BBP shows an error message saying "Path not found" and await for the user to return back or changing query \\ \hline
\end{tabularx}
\caption{Use Case 9: Search Path}
\label{tab:uc9}
\end{table}

\begin{figure}[H]
    \centering
    % \includegraphics[width=\textwidth]{uc9_sequence.png}
    \caption{Sequence Diagram for UC9: Search Path}
    \label{fig:uc9_sequence}
\end{figure}

\subsubsection{[UC10] - Visualize Path}

\begin{table}[H]
\centering
\begin{tabularx}{\textwidth}{|l|X|}
\hline
\textbf{Participating Actors} & Any User (Guest/Registered), Map Service, Weather Service \\ \hline
\textbf{Entry Conditions} & The User has searched paths between two points and has chosen one listed option \\ \hline
\textbf{Flow of events} &
\begin{enumerate}
    \item The user clicks on one path in the search results list
    \item BBP fetches map data about chosen path's coordinates
    \item BBP queries external Weather Service about current conditions along the chosen path
    \item BBP displays a detailed view showing path trajectory on the map, aggregated statistics (average time, total distance and status), confirmed potholes and current weather conditions along the path
\end{enumerate} \\ \hline
\textbf{Exit Conditions} & User is shown of path on the map, its aggregated statistics and current weather conditions \\ \hline
\textbf{Exceptions} &
Weather Service is not available: BBP shows an error regarding weather conditions but still loads everything else

Map Service is unavailable: BBP shows an error regarding the map but still shows all other things \\ \hline
\end{tabularx}
\caption{Use Case 10: Visualize Path}
\label{tab:uc10}
\end{table}

\begin{figure}[H]
    \centering
    % \includegraphics[width=\textwidth]{uc10_sequence.png}
    \caption{Sequence Diagram for UC10: Visualize Path}
    \label{fig:uc10_sequence}
\end{figure}

\section{Performance Requirements}
\label{sec:performance}

To ensure an optimal user experience, the system should complete actions within 2 seconds. However for this application, that does not involve critical services, strict requirements are not needed. For this reason completion time for complex tasks can be extended up to 5 seconds. Those complex tasks could be:

\begin{itemize}
    \item searching and displaying available paths
    \item calculating and displaying statistical information
\end{itemize}

\section{Design Constraints}
\label{sec:constraints}

\subsection{Standards Compliance}

Since the system collects and stores sensitive user data, the design must strictly adhere to the General Data Protection Regulation (GDPR). This includes implementing mechanisms for:
\begin{itemize}
    \item Explicit user consent for tracking.
    \item The deletion of all user data upon request.
    \item Anonymization of data when aggregated for public statistics.
\end{itemize}

\subsection{Hardware Limitations}

\begin{itemize}
    \item \textbf{Mandatory Sensor Availability}: The automated tracking module is constrained by the physical hardware of the user's device. The application cannot function in automated mode on devices lacking working sensors such as GPS, accelerometer, gyroscope.

    \item \textbf{Local Storage Capacity}: The local buffering feature for offline recording is constrained by the available internal storage of the mobile device. The application design must include a safe upper limit for the buffer size to avoid filling the user's device storage.

    \item \textbf{Battery Drain Constraints}: The software must use optimizations techniques or/and APIs to minimize battery usage.
\end{itemize}

\subsection{Implementation Constraints}

\begin{itemize}
    \item \textbf{Third-party APIs Limits}: The design is constrained by the Rate Limits imposed by external service providers. The software must implement some strategies to avoid exceeding Rate Limits and incurring service denials or extra costs.
\end{itemize}

\section{Software System Attributes}
\label{sec:attributes}

\subsection{Reliability}

To ensure better reliability performance, regular maintenance and bug fixes should be performed. Planned maintenance should also be notified to users to avoid any unexpected inconvenience.

\subsection{Availability}

To ensure a positive user experience, the system uptime should be 99\% or more. This application does not involve critical tasks (e.g., financial transactions) so the minimum standard it's not required to be 99.5\% or more.

The core function of trip recording in automated mode should be available 100\% of the time, regardless of network status, provided the device has sufficient battery and storage.

\subsection{Security}

The system shall enforce strict access control policies. A user's specific history and sensor logs must be accessible only to that specific user, unless explicitly published.

User passwords shall never be stored in plain text. The system must use industry-standard hashing algorithms with unique salts for credential storage.

\subsection{Maintainability}

The architecture should follow a modular approach to facilitate debugging and programming.

The source code should also be documented following standard conventions (e.g., Javadoc for Java, PyDoc for Python) to facilitate future maintenance and onboarding of new developers.

\subsection{Portability}

BBP is a mobile application so it has to be compatible with both major OSs (Android and iOS).
